\section{Conclusiones}

El diseño y prototipo de la aplicación VITA demuestran que la aplicación rigurosa de metodologías de diseño centrado en el humano es fundamental para eliminar la fricción técnica \parencite{diaz2013diseno,iso9241_210}.
La implementación de atajos visuales, selectores interactivos y validaciones restrictivas en tiempo real no solo facilita la curva de aprendizaje, sino que asegura la recolección de datos puros que alimentan las proyecciones metabólicas del usuario. Al integrar principios de accesibilidad y componentes de gamificación, la aplicación garantiza una alta satisfacción, cumpliendo exitosamente con los estándares de calidad para promover la constancia en la actividad física \parencite{iso9241_11,mejia2024calidad}.